% Source: https://artofproblemsolving.com/community/c3h2217460_an_probability_problem
% Stable identity: aops-topic-2217460-q01
\textbf{One rainy day.} Let $T$ be the number of days up to and including the first rainy day, and $E=\mathbb{E}[T]$. On the first day, it rains with probability $\dfrac{1}{5}$ (then $T=1$); otherwise, with probability $\dfrac{4}{5}$, one day has passed and the situation starts afresh. So
\[E=\dfrac{1}{5}\cdot1+\dfrac{4}{5}(1+E)\quad\Longrightarrow\quad\dfrac{1}{5}E=1\quad\Longrightarrow\quad E=5.\]

\textbf{Thirty rainy days.} Split the waiting time into $30$ stretches: from the start to the first rainy day, from the day after the first rainy day to the second, and so on. By independence, each stretch has the same distribution as $T$, so each has expected length $5$. By linearity of expectation the total expected time is
\[30\cdot5=\boxed{150}\text{ days}.\]

\textbf{Intuition.} Rain falls on one day in five on average, so to collect $30$ rainy days we expect to need $30\cdot5$ days. (This is an expectation: in a particular year it may take more or fewer days.)
